\subsection{Guidance-interface contract and transferability boundary}
\label{subsec:guidance_interface_transferability}

The reported numerical results are obtained with the frozen success-latched guidance implementation used in this study. However, the dwell-confirmed validation layer is not tied to the internal phase labels or terminal-capture variables of that implementation. To make this separation explicit, the validation layer can be expressed as an interface contract between an arbitrary guidance algorithm and the post-processed TAEM closure checker.

The required interface is the logged time sequence, a vehicle identifier, and the TAEM evaluation channels needed to reconstruct altitude, velocity, and range-to-go errors relative to the same target box. Given these fields, the checker evaluates the in-box indicator, updates the \(N_{\mathrm{dwell}}\)-sample dwell counter, latches the first accepted terminal event, and applies the terminal-reason and close-pass-escape checks when those fields are available. This procedure does not require the guidance algorithm to use the same internal phases, terminal-capture law, bank-command smoothing, or route-bias implementation.

The interface audit in Table~\ref{tab:transferability_interface_audit} replayed the final sampled-validation logs using only guidance-agnostic terminal-state interface fields. The portable replay reconstructed 149/150 = 0.993333 vehicle-level dwell closures and 49/50 = 0.980000 case-level dwell closures. The strict-like replay, which additionally uses available terminal-reason and close-pass-escape fields, reconstructed 149/150 = 0.993333 vehicle-level closures and 49/50 = 0.980000 case-level closures. No V8.2-specific phase or terminal-capture field was consumed by the portable replay.

The transferability claim is therefore limited but explicit. The core validation ideas---terminal-state box membership, sample-count dwell, event latching, staged evidence separation, and boundary-diagnostic reporting---can be applied to any guidance algorithm that exports the required terminal-state interface. The numerical success rates reported in this paper, however, remain specific to the frozen guidance implementation, scenario, tolerance box, and sampled local support used here. A different guidance law would have to be re-run and re-audited under the same interface contract before any performance claim could be transferred.
